\slajdid{04-s026}
\begin{frame}[label=04-s026]{Gde je binarna tačka?}
% element: 04-s026:e01 — Pitanje binarne tačke i mogućnost celobrojnog, razlomljenog ili mešovitog zapisa
% element: 04-s026:e02 — Registar 01101000 i nemogućnost posebnog koda tačke sa 0 i 1
% element: 04-s026:e03 — Projektant predviđa položaj; oba podatka o značenju moraju biti poznata
% element: 04-s026:e04 — Primeri tačke iza D5, iza D0, ispred D7 i drugi položaji
Registar može sadržati celobrojnu, razlomljenu ili obe vrednosti.
\par\vspace{1.5mm}{\fontsize{10}{12}\selectfont\renewcommand{\arraystretch}{.95}\begin{tabularx}{\linewidth}{*{8}{>{\centering\arraybackslash}X}}\toprule
$D_7$&$D_6$&$D_5$&$D_4$&$D_3$&$D_2$&$D_1$&$D_0$\\\midrule
0&1&1&0&1&0&0&0\\\bottomrule\end{tabularx}
\par}\par\vspace{1.5mm}
Raspolažemo samo sa $0$ i $1$; tačka se tim ciframa ne može označiti.
\Rezultat{Položaj binarne tačke određuje projektant; mora biti poznat.}
\par\vspace{1.5mm}
{\fontsize{10}{12}\selectfont\renewcommand{\arraystretch}{.95}\begin{tabularx}{\linewidth}{@{}lX@{}}\toprule
Položaj tačke & Značenje istog registra\\\midrule
Iza $D_5$ & $011.01000_2$\\
Iza $D_0$ & Celobrojna vrednost\\
Ispred $D_7$ & Razlomljena vrednost\\
\ldots & Drugi predviđeni položaj\\\bottomrule
\end{tabularx}\par}
\end{frame}
